\chapter{Vulnerability Overview}

\section{Selection Criteria}

From the prior team's assessment of 60+ potential vulnerabilities, five were selected based on:
\begin{enumerate}[label=\arabic*.]
  \item \textbf{High verifier bypass probability}: Previously classified as "exploitable" or "limited"
  \item \textbf{Representativeness}: Cover multiple vulnerability classes (ABI, type confusion, bounds checking, taint tracking)
  \item \textbf{Practical feasibility}: Can be triggered via network packets within the test infrastructure
  \item \textbf{Evidence potential}: Clear stages from bytecode to runtime where impact is observable
\end{enumerate}

\section{Vulnerability Summary Table}

\noindent
\begin{center}
\begin{tabular}{|c|l|l|l|}
\hline
\textbf{ID} & \textbf{Vulnerability Name} & \textbf{ISO Rule} & \textbf{Class} \\
\hline
5.06a & Function Pointer Type Mismatch & 5.06 & ABI (Application Binary Interface) Mismatch \\
5.06b & Wrong Argument Count & 5.06 & ABI (Application Binary Interface) Mismatch \\
5.06d & Wrong Argument Types & 5.06 & ABI (Application Binary Interface) Mismatch \\
5.39 & Unprototyped Function Pointer & 5.39 & Type Confusion (Stack OOB Write) \\
5.46b & Tainted Length (Packet Buffer) & 5.46 & Bounds Checking (Packet OOB Write) \\
\hline
\end{tabular}
\end{center}


\section{Vulnerability Classes}

\subsection{ABI Mismatch Vulnerabilities (5.06a, 5.06b, 5.06d)}

These three vulnerabilities exploit discrepancies between function signatures and actual function call invocations.

\subsubsection{ISO-IEC Rule 5.06 Background}

Rule 5.06 specifies: \textit{``Do not declare pointers to functions with a type that differs from the function definition.''}

Violating this rule means a function pointer is declared with one signature (e.g., \texttt{int (*fp)(int)}) while the actual function has a different signature (e.g., \texttt{int func(int, int, int)}). When calling through the pointer, mismatched numbers and types of arguments are passed, violating ABI expectations.

\subsubsection{Why This Matters for eBPF}

The eBPF verifier must validate that:
\begin{enumerate}[label=\arabic*.]
  \item Function pointers are declared with correct types
  \item Calls through pointers pass arguments of the correct count and type
  \item Registers align correctly with the expected function signature (x86-64 ABI (Application Binary Interface))
\end{enumerate}

If the verifier fails to catch signature mismatches, it may allow registers to be used as the wrong type (e.g., uninitialized register passed as pointer), extra arguments to overwrite adjacent stack or registers, or type truncation or sign-extension bugs (e.g., 64-bit pointer passed as 32-bit int).

\subsubsection{Classification of 5.06a, 5.06b, 5.06d}

Based on prior analysis and practical testing, vulnerability 5.06a represents type incompatibility where the function expects a different return type, 5.06b involves wrong argument counts where calls are made with more arguments than declared, and 5.06d involves wrong argument types such as passing an int where a long is expected.

All three are \textbf{runtime triggers} — the verifier accepts the code despite the mismatch, and runtime behavior is observable via kernel trace output. However, they do not constitute true \textbf{memory corruption} in the test environment (the vulnerability is in register alignment, not buffer overflow).

\subsection{Type Confusion with Stack OOB Write (5.39 taintnoproto)}

\subsubsection{ISO-IEC Rule 5.39 Background}

Rule 5.39: \textit{``Do not create pointers to functions with an implicit return type.''}

In this violation, a function pointer is declared without a complete prototype (missing argument types). The pointer is called with specific arguments, but the compiler and verifier cannot properly validate the call because the prototype is incomplete.

\subsubsection{The 5.39 Exploit Mechanism}

The XDP SynProxy vulnerability (5.39) combines:
\begin{enumerate}[label=\arabic*.]
  \item A TCP sequence number (64-bit, attacker-controlled via SYN packet)
  \item An array of \texttt{u16} values on the stack
  \item An unprototyped function pointer call that uses the sequence number as an array index
\end{enumerate}

If the sequence number is large (e.g., 18, 25, etc.), the array access goes out-of-bounds and writes to adjacent stack memory.

\subsubsection{Classification}

\textbf{5.39 is an EXPLOITABLE vulnerability} (not just a runtime trigger). The \textbf{out-of-bounds write} targets the eBPF program's local stack frame, allowing an attacker who controls the sequence number to determine which eBPF-local stack location is overwritten. However, due to the eBPF sandbox model, such corruption is limited to the XDP handler's execution context and cannot directly compromise kernel memory or enable privilege escalation. The vulnerability is exploitable for Denial of Service or packet processing disruption.

\subsection{Bounds Checking with Tainted Length (5.46b taintsink)}

\subsubsection{ISO-IEC Rule 5.46 Background}

Rule 5.46: \textit{``Do not apply the sizeof operator to a polymorphic object.''}

More broadly, rule 5.46-adjacent violations include using user-controlled (tainted) values for memory copy lengths (\texttt{memcpy(dst, src, user\_len)}), array indices, and loop bounds.

\subsubsection{The 5.46b Exploit Mechanism}

The XDP SynProxy vulnerability (5.46b) is:
\begin{enumerate}[label=\arabic*.]
  \item Parsing a TCP header payload length from an incoming packet
  \item Using that length directly in a loop that writes to packet buffer fields
  \item The destination field is fixed-size (e.g., 6-byte MAC address field)
  \item An attacker crafts a packet with a large payload length (e.g., 60 bytes from TCP doff field)
  \item Result: Packet buffer overflow, overwriting adjacent packet header fields (MAC source, protocol)
\end{enumerate}

\subsubsection{Classification}

\textbf{5.46b is an EXPLOITABLE vulnerability}: The \textbf{packet-length-based buffer overflow} is directly triggered and controlled by an attacker. The verifier accepts bounds checks on tainted values, allowing the write size to exceed the buffer. Attackers can overflow 54+ bytes past the 6-byte \texttt{h\_dest} field, corrupting adjacent \textbf{packet buffer memory} (within the same packet structure). This overflow is confined to in-kernel packet buffer operations and does NOT directly compromise kernel data structures or enable privilege escalation. The vulnerability demonstrates a verifier gap but has limited real-world impact.
\begin{itemize}
  \item Direct packet buffer overflow
  \item Attacker controls the overflow size (via TCP doff field)
  \item Overwrites adjacent packet fields (MAC source, protocol)
  \item Confirmed memory write primitive within packet scope only
\end{itemize}

\section{Evidence and Severity Breakdown}

\noindent
\begin{center}
\begin{tabular}{|c|l|c|l|}
\hline
\textbf{ID} & \textbf{Name} & \textbf{Severity} & \textbf{Evidence Quality} \\
\hline
5.06a & Type mismatch & Medium & Bytecode + Verifier + Runtime \\
5.06b & Wrong arg count & Medium & Bytecode + Verifier + Runtime \\
5.06d & Wrong arg types & Medium & Bytecode + Verifier + Runtime \\
5.39 & Stack OOB (eBPF-local) & Medium-High & Bytecode + Verifier + Runtime (OOB write confirmed) \\
5.46b & Packet Buffer OOB & Medium & Bytecode + Verifier + Runtime (Overflow confirmed) \\
\hline
\end{tabular}

\textbf{Table 6.2: Severity and evidence quality for each vulnerability}
\end{center}

\section{Exploitation Chains}

\subsection{5.06a/5.06b/5.06d: Potential Escalation Path}

While these are runtime triggers in the current target, they could become exploitable in a chained attack:
\begin{enumerate}[label=\arabic*.]
  \item Trigger the ABI mismatch to observe register/stack state
  \item Use information disclosure to leak kernel addresses or structure offsets
  \item Combine with a separate vulnerability for full exploitation
\end{enumerate}

\subsection{5.39: Direct Exploitation}

Stack OOB write is confirmed as a memory corruption primitive within the eBPF sandbox. Full exploitation for DoS would require:
\begin{enumerate}[label=\arabic*.]
  \item Knowledge of eBPF stack frame structure
  \item Identification of exploitable eBPF-local data structures
  \item Triggering specific sequence numbers to corrupt the target
\end{enumerate}

\subsection{5.46b: Direct Exploitation}

Packet buffer overflow provides a confirmed memory write primitive within packet scope:
\begin{enumerate}[label=\arabic*.]
  \item Attacker controls overflow size via crafted TCP Data Offset field
  \item Attacker controls overflow data via packet payload contents
  \item Corrupts in-kernel packet buffer memory (adjacent header fields)
  \item With XDP processing of attacker-controlled packets, the vulnerability is reliably triggerable
  \item \textbf{NOTE}: Overflow is limited to packet-local memory; does NOT enable kernel-level privilege escalation
\end{enumerate}

\section{Impact Differentiation: Runtime Triggers vs.\ Memory Corruption Primitives}

A critical distinction emerges when comparing the five vulnerabilities:

\subsection{Runtime Triggers (Limited Impact): 5.06a, 5.06b, 5.06d}

The three ABI mismatch vulnerabilities are \textbf{verifier bypasses} but have \textbf{limited practical impact} in the XDP SynProxy context:

\begin{itemize}
  \item \textbf{No direct memory corruption}: The mismatched calls do not cause buffer overflows or out-of-bounds writes.
  \item \textbf{No control flow hijacking}: Function pointers still target real, kernel-intended functions (not attacker addresses).
  \item \textbf{Trigger-able but not exploitable}: The vulnerabilities can be triggered via network packets and observed via kernel traces, but they do not yield a usable memory write primitive.
  \item \textbf{Contextual threat}: If similar patterns appeared in code using the mismatched arguments for pointer arithmetic, bounds calculations, or array indexing, they would become more dangerous. In this PoC, they are demonstrations of verifier design gaps rather than exploitable vulnerabilities.
\end{itemize}

\subsection{Memory Corruption Primitives (Significant Impact): 5.39, 5.46b}

The two taint-tracking vulnerabilities are \textbf{confirmed exploitable} and provide \textbf{memory corruption capability} with \textbf{scope limitations}:

\begin{itemize}
  \item \textbf{5.39 (Stack OOB Write)}: Demonstrated write to out-of-bounds eBPF-local stack locations with attacker-controlled offset (derived from TCP sequence number). \textbf{Limited to eBPF sandbox}; cannot compromise kernel memory.
  \item \textbf{5.46b (Packet Buffer Overflow)}: Demonstrated 54-byte overflow beyond a 6-byte field with attacker-controlled data (packet payload). \textbf{Limited to packet-local memory} (adjacent packet fields); does not enable kernel-level RCE (Remote Code Execution).
  \item \textbf{Impact within scope}: 5.39 can disrupt XDP packet processing; 5.46b can corrupt packet headers in flight.
  \item \textbf{Reliability}: Both are reliably triggerable via standard network packets; no special conditions are required.
\end{itemize}

\subsection{Summary}

\begin{center}
\begin{tabular}{|c|c|l|}
\hline
\textbf{Vulnerability} & \textbf{Type} & \textbf{Practical Impact in SynProxy} \\
\hline
5.06a, 5.06b, 5.06d & Runtime Trigger & Verifier bypass; limited exploitability \\
5.39 & Memory Corruption & Stack OOB write primitive; highly exploitable \\
5.46b & Memory Corruption & Buffer overflow primitive; highly exploitable \\
\hline
\end{tabular}
\end{center}

\section{Next Steps}

The following chapters provide detailed analysis of each vulnerability, including:
\begin{itemize}
  \item Complete source code of the vulnerability patch
  \item Bytecode and verifier analysis showing that the issue is preserved
  \item PoC code and runtime evidence
  \item Assessment of real-world exploitability
\end{itemize}
